\begin{remark}[5.11.3]
\label{IV.5.11.3}
The conclusion that $A^{(1)}$ is a finite $A$-algebra need no
longer hold if, in hypothesis {\rm(ii)} of
\sref{IV.5.11.2}, one assumes only that $A$ is catenary.

Indeed, consider the catenary local ring $A$ constructed in
\sref{IV.5.6.11}; its integral closure $A'$ (denoted there by
$B$) is a finite $A$-algebra.  If $A^{(1)}$ were also finite over
$A$, then, being contained in the fraction field of $A$, it would
be contained in $A'$.  On the other hand, in the notation of
\sref{IV.5.6.11}, every height-one prime ideal of $A$ is of the
form $\mathfrak pA$, where $\mathfrak p$ is a height-one prime of
$C$, and $A_{\mathfrak pA}=C_{\mathfrak p}$.
Moreover, $C_{\mathfrak p}=E_{\mathfrak p'}$, where
$\mathfrak p'$ is the unique prime of $E$ above $\mathfrak p$.
Thus $A_{\mathfrak pA}$ is integrally closed and consequently
contains $A'$.  By definition this gives
$A'\subset A^{(1)}$, so finiteness of $A^{(1)}$ would force
$A^{(1)}=A'$.  This is impossible: one of the two prime ideals of
$A'$ above the maximal ideal $\mathfrak nA$ of $A$ has height
$1$, whereas $\mathfrak nA$ has height $2$, contradicting
\sref{IV.5.10.17},~(iv).
\end{remark}

\begin{corollary}[5.11.4]
\label{IV.5.11.4}
Let $X$ be a locally Noetherian prescheme, $Z$ a closed subset of $X$,
$U=X\setmin Z$, $i:U\to X$ the canonical injection, and $\sh F$ a
coherent $\sh O_U$-Module.  For $i_*(\sh F)$ to be a coherent
$\sh O_X$-Module, it is necessary that, for every
$x\in\operatorname{Ass}(\sh F)$, we have
\[
  \codim\bigl(\overline{\{x\}}\cap Z,\overline{\{x\}}\bigr)
  \geq2.
\]
This condition is sufficient in each of the following two cases:
\begin{enumerate}[{\rm(i)}]
  \item The prescheme $X$ is locally immersible in a regular scheme.
  \item The prescheme $X$ is universally catenary and universally
  Japanese \textup{(which means, by definition, that every point
  $x\in X$ has an affine open neighbourhood whose ring is}
  universally Japanese \textup{\sref[0]{0.23.1.1})}.
\end{enumerate}
\end{corollary}

\begin{proof}
We know \sref[I]{I.9.4.7} that there is a \emph{coherent}
sub-$\sh O_X$-Module $\sh G$ of $i_*(\sh F)$ such that
$\sh G|U=\sh F$.  We evidently have $\operatorname{Ass}(\sh F)
\subset\operatorname{Ass}(\sh G)
\subset\operatorname{Ass}\bigl(i_*(\sh F)\bigr)$,
and since
\[
  \operatorname{Ass}\bigl(i_*(\sh F)\bigr)
  =\operatorname{Ass}(\sh F)
\]
\sref{IV.3.1.13}, we have
$\operatorname{Ass}(\sh G)=\operatorname{Ass}(\sh F)$;
it then suffices to apply \sref{IV.5.11.1} to the coherent
$\sh O_X$-Module $\sh G$, noting that $i_*(\sh F)=\sh H^0_{X/Z}(\sh G)$
and that, when $X$ is universally catenary and universally
Japanese, hypothesis {\rm(ii)} of \sref{IV.5.11.1} is satisfied by
definition.
\end{proof}

\begin{corollary}[5.11.5]
\label{IV.5.11.5}
Let $X$ be a locally Noetherian prescheme, $Z$ a subset
of $X$ stable under specialization, and $\sh F$ a coherent
$\sh O_X$-Module such that
$\operatorname{Ass}(\sh F)\subset X\setmin Z$.  Then the condition:
\begin{enumerate}[{\rm a)}]
  \item $\sh H^0_{X/Z}(\sh F)$ is a coherent
  $\sh O_X$-Module
\end{enumerate}
implies the following:
\begin{enumerate}
  \item[{\rm d)}] For every subset $T$ of $Z$ closed in $X$
  \textup{(or only for a filtered}

\oldpage[IV-2]{125}

  increasing family \textup{of closed subsets $T$ of $Z$, with union
  $Z$)}, $i_*(\sh F|X\setmin T)$ \textup{(}where
  $i:X\setmin T\to X$ is the canonical injection\textup{)} is coherent.
\end{enumerate}
When $X$ satisfies one of the hypotheses \textup{(i), (ii)} of
\sref{IV.5.11.1}, \textup{a)} and \textup{d)} are equivalent.
\end{corollary}

\begin{proof}
Note that we have
\[
  \operatorname{Ass}\bigl(i_*(\sh F|X\setmin T)\bigr)
  =\operatorname{Ass}(\sh F)
\]
by the hypothesis and \sref{IV.3.1.13}; it therefore follows from
\sref{IV.5.10.2} that the canonical maps
\[
  i_*(\sh F|X\setmin T)\longrightarrow\sh H^0_{X/Z}(\sh F)
\]
are \emph{injective}; the fact that \emph{a)} implies \emph{d)} is
therefore a consequence of this remark.  Conversely, condition
\emph{d)} implies, by \sref{IV.5.11.1}, that
$\codim(T\cap Y_\alpha,Y_\alpha)\geq2$ with the notation of
\sref{IV.5.11.1}; consequently
$\codim(Z\cap Y_\alpha,Y_\alpha)\geq2$, since $Z$ is the union of
its subsets which are closed in $X$, and the last assertion of the
corollary follows from \sref{IV.5.11.1}.
\end{proof}

\begin{corollary}[5.11.6]
\label{IV.5.11.6}
Let $A$ be a Noetherian ring and $X=\Spec(A)$.  Consider the
following properties:
\begin{enumerate}[{\rm a)}]
  \item For every integral quotient ring $B$ of $A$, the ring $B^{(1)}$
  \textup{(notation of \sref{IV.5.11.2})} is a finite $B$-algebra.
  \item For every coherent $\sh O_X$-Module $\sh F$ and every
  subset $Z$ of $X$, stable under specialization, and such that for
  every $x\in\operatorname{Ass}(\sh F)\cap(X\setmin Z)$ we have
  $\codim\bigl(\overline{\{x\}}\cap Z,\overline{\{x\}}\bigr)\geq2$,
  the $\sh O_X$-Module $\sh H^0_{X/Z}(\sh F)$ is coherent.
  \item For every closed subset $T$ of $X$ and every coherent
  $\sh O_U$-Module $\sh G$ \textup{(}where $U=X\setmin T$\textup{)} such
  that for every $x\in\operatorname{Ass}(\sh G)$ we have
  $\codim\bigl(\overline{\{x\}}\cap T,\overline{\{x\}}\bigr)\geq2$,
  $i_*(\sh G)$ \textup{(}where $i:U\to X$ is the canonical
  injection\textup{)} is a coherent $\sh O_X$-Module.
  \item For every integral quotient ring $B$ of $A$ and every ideal
  $\mathfrak J$ of height $\geq2$ in $B$, the ring
  $\bigcap_{\mathfrak p\not\supset\mathfrak J}B_{\mathfrak p}$
  is a finite $B$-algebra.
\end{enumerate}
We then have the implications
\[
  {\rm a)}\Longleftrightarrow{\rm b)}
  \Longrightarrow{\rm c)}
  \Longleftrightarrow{\rm d)}.
\]
Moreover, conditions \textup{a), b), c)} and \textup{d)} are satisfied
in each of the following two cases:
\begin{enumerate}[{\rm(i)}]
  \item $A$ is a quotient of a regular ring.
  \item $A$ is universally catenary and universally Japanese.
\end{enumerate}
\end{corollary}

\begin{proof}
Let $B=A/\mathfrak q$, where $\mathfrak q$ is a prime ideal of
$A$, so that $Y=\Spec(B)$ is the closed subset $V(\mathfrak q)$ of
$X$; put $Z=Z^{(2)}(Y)$, which is a subset of $X$ stable under
specialization; since $B$ is integral,
$\operatorname{Ass}(A/\mathfrak q)$ reduces to the generic point
$\mathfrak q$ of $Y$.  If condition \emph{b)} is satisfied, we can
apply it to the coherent $\sh O_X$-Module $\sh F=(A/\mathfrak q)^{\sim}$
and to $Z$, and by \sref{IV.5.9.3.1} this shows that $B^{(1)}$
is an $A$-module of finite type, and \emph{a fortiori} a $B$-module of
finite type.  Conversely, suppose \emph{a)} satisfied; then, if $\sh F$
is a coherent $\sh O_X$-Module such that for every
$x\in\operatorname{Ass}(\sh F)\cap(X\setmin Z)$ we have
$\codim\bigl(\overline{\{x\}}\cap Z,\overline{\{x\}}\bigr)\geq2$, we can
apply (with the notation of \sref{IV.5.11.1}) the result of \emph{a)}
to each of the affine schemes $Y_\alpha=\Spec(B_\alpha)$, where
$B_\alpha$ is an integral quotient ring of $A$; since by hypothesis
$Z_\alpha$ is contained in $Z^{(2)}(Y_\alpha)$, condition \emph{a)}
\textup{(}taking into account \sref{IV.5.10.2} and the fact that
$\operatorname{Ass}(\sh O_{Y_\alpha})$ reduces to the generic point of
$Y_\alpha$\textup{)} implies that
$\sh H^0_{Y_\alpha/Z_\alpha}(\sh O_{Y_\alpha})$ is a coherent
$\sh O_{Y_\alpha}$-Module, and \emph{b)} then follows from
\sref{IV.5.11.1}.

\oldpage[IV-2]{126}

To see that \emph{c)} implies \emph{d)}, we reason as above, applying
\emph{c)} to the case where $\sh G=(A/\mathfrak q)^{\sim}|U$ and
$T=V(\mathfrak J)$; conversely, one proves that \emph{d)} implies
\emph{c)} by again using the equivalence of \emph{a)} and \emph{b)} in
\sref{IV.5.11.1}.  It is evident that \emph{c)} is a particular case of
\emph{b)}.  Finally, the fact that \emph{a)} \textup{(}and consequently
each of the other conditions\textup{)} is satisfied when $A$ satisfies one
of the hypotheses {\rm(i)}, {\rm(ii)} follows from \sref{IV.5.11.2}, in
view of the definition of universally catenary rings and of universally
Japanese rings.
\end{proof}

\begin{remarks}[5.11.7]
\label{IV.5.11.7}
\begin{enumerate}[{\rm(i)}]
  \item It is not known whether, in \sref{IV.5.11.5}, condition
  \emph{d)} implies \emph{a)} without an additional hypothesis on $X$;
  we shall see later \sref{IV.7.2.4} that this is indeed so when $X$
  is a \emph{local} scheme.  Likewise, we shall prove that when $A$ is
  a Noetherian \emph{local} ring the four properties \emph{a)},
  \emph{b)}, \emph{c)}, \emph{d)} of \sref{IV.5.11.6} are equivalent
  \sref{IV.7.2.4}.  It is not known whether this result extends to all
  Noetherian rings.

  \item If $A$ satisfies property \emph{a)} of
  \sref{IV.5.11.6}, so does every \emph{ring of fractions} $S^{-1}A$
  and every \emph{finite} $A$-algebra $C$.  Indeed, every integral
  quotient ring of $S^{-1}A$ is of the form $S^{-1}(A/\mathfrak q)$,
  where $\mathfrak q$ is a prime ideal of $A$ not meeting $S$; and on
  the other hand, if $\mathfrak r$ is a prime ideal of $C$ and
  $\mathfrak p$ its inverse image in $A$, $C/\mathfrak r$ is a finite
  integral $(A/\mathfrak p)$-algebra containing $A/\mathfrak p$; our
  assertions are therefore consequences of
  \sref{IV.5.10.17},~(iii) and~(iv).
\end{enumerate}
\end{remarks}
